\chapter{Domain and Language Adaptation}
\label{ch:domain-language-adaptation}

\abstract{This chapter studies how LLMs are adapted to Chinese, domain-specific corpora, NL2SQL, medical-style retrieval tasks, and other specialized settings. It compares continual pretraining, supervised fine-tuning, retrieval, tokenizer extension, and evaluation under domain shift.}

\section{Chapter Spine}
Domain adaptation is not one technique. The right choice depends on whether the missing capability is vocabulary coverage, factual knowledge, style, reasoning format, schema grounding, or access to private documents.

\section{Core Sections}
\subsection{Language Adaptation}
Tokenizer extension, bilingual corpora, Chinese LLaMA-style workflows, and multilingual evaluation.

\subsection{Domain Knowledge}
Continual pretraining versus retrieval versus SFT for specialized knowledge.

\subsection{Structured Tasks}
NL2SQL, schema linking, execution accuracy, and dataset leakage.

\subsection{Safety in Domains}
Medical, legal, financial, and organizational settings as high-stakes contexts requiring stronger evaluation.

\section{Planned Exercises}
Compare a retrieval-only baseline with a LoRA-adapted model on a small domain QA set.
